\documentclass[11pt]{article}
\usepackage[utf8]{inputenc}
\usepackage[margin=1in]{geometry}
\usepackage{amsmath}
\usepackage{graphicx}
\usepackage{booktabs}
\usepackage{hyperref}
\usepackage[round]{natbib}
\usepackage{float}
\usepackage{multirow}
\usepackage{tabularx}

\title{Stable Population Structure in Europe Since the Iron Age, Despite High Mobility}
\author{}
\date{}

\begin{document}
\maketitle

\begin{abstract}
Ancient DNA studies have revealed dramatic population turnovers during European prehistory, including the Neolithic farmer expansion and Bronze Age Steppe migrations. However, the dynamics of population structure during the historical period (the last $\sim$3,000 years) remain poorly understood despite evidence for extensive individual mobility. Here, we report 204 new whole genomes (median depth 0.92$\times$; range 0.16--2.38$\times$) from 53 archaeological sites across 18 countries spanning the Iron Age through the post-medieval period, including the first historical-period genomes from Armenia, Algeria, Austria, and France. By analyzing these alongside 2,033 present-day, 1,998 prehistoric, and 764 published historical genomes ($n \approx 5{,}000$ total), we demonstrate that the overall geographic structure of European genetic variation has remained remarkably stable from the Iron Age to the present. Yet at least 7\% of historical individuals carry ancestry uncommon in their sampling region, indicating extensive cross-regional mobility, particularly during Imperial Rome. Stepping-stone migration simulations show that this level of apparent dispersal ($\geq$7\% per generation under panmictic assumptions) should rapidly collapse population structure ($p < 0.001$, permutation test), implying that effective migration was substantially lower than apparent dispersal. We propose that improved transportation networks enabled transient movement of individuals---for trade, military service, and labor---without proportional gene flow, reconciling high individual mobility with population-level genetic stability.
\end{abstract}

\section{Introduction}

The population genetic landscape of modern Europe was shaped by a series of major demographic events during prehistory. Three ancestral components---Western Hunter-Gatherers (WHG), Neolithic Farmers originating from Anatolia around 7,500 BCE, and Bronze Age Steppe Herders arriving approximately 3,500 BCE---admixed to form the tripartite ancestry of present-day Europeans \citep{lazaridis2014ancient, haak2015massive, allentoft2015population, mathieson2015genome}. These components were largely established by the end of the Bronze Age (circa 1,000 BCE) and their geographic distribution tracks modern population structure \citep{fu2016genetic, lazaridis2016genomic, posth2023paleogenomics}.

While the prehistoric period has been densely sampled by ancient DNA studies, the historical period---spanning roughly the last 3,000 years---remains comparatively understudied from a pan-European perspective. Individual regional studies have documented high genetic heterogeneity during the Roman Imperial period in Italy \citep{antonio2019ancient}, the Iberian Peninsula \citep{olalde2019genomic}, and the central Mediterranean \citep{moots2023genetic}. Viking-era genomes from Scandinavia revealed widespread mobility across northern Europe \citep{margaryan2020population}. However, no comprehensive multi-region assessment has addressed how overall population structure changed (or persisted) across western Eurasia during historical times.

This gap is especially important because the historical period brought unprecedented connectivity. The Roman Empire, at its height, connected territories from Britain to Mesopotamia through an extensive road and maritime network. Travel times that spanned centuries during prehistoric migrations could be completed in weeks \citep{scheidel2019escape, orbis}. If even a modest fraction of this mobility translated into gene flow, standard population genetics models predict rapid erosion of geographic structure \citep{kimura1964stepping, wright1943isolation}. Yet modern Europeans display clear geographic genetic structure that closely mirrors the map of Europe \citep{price2006principal}, suggesting that either historical mobility was genetically inconsequential or that apparent dispersal overestimates effective migration.

To resolve this paradox, we generated 204 new whole genomes from historical-period individuals across 18 countries and analyzed them within a dataset of approximately 5,000 genomes spanning the Mesolithic to the present. We assessed population structure using PCA projection, ancestry decomposition via qpAdm, and systematic outlier detection, then tested whether observed dispersal rates are compatible with structural stability using stepping-stone migration simulations.

\section{Results}

\subsection{Genome-Wide Data from 204 Historical Individuals}

We generated whole-genome shotgun data from 204 individuals excavated at 53 archaeological sites across 18 countries (Table~\ref{tab:sample_overview}). DNA was extracted primarily from petrous bones ($n = 203$) with one extraction from teeth, using partial UDG treatment to reduce deamination artifacts while preserving limited deamination patterns for authentication \citep{rohland2015partial}. Pseudohaploid genotypes were called on the 1240k SNP panel, yielding a median of 685,058 SNPs per sample (range: 167,000--1,029,345). Radiocarbon dating was performed on 126 samples. These genomes include the first historical-period ancient DNA from Armenia, Algeria, Austria, and France.

\begin{table}[H]
\centering
\caption{Sample overview and sequencing statistics for 204 newly reported historical-period genomes.}
\label{tab:sample_overview}
\begin{tabular}{lc}
\toprule
\textbf{Parameter} & \textbf{Value} \\
\midrule
New genomes reported & 204 \\
Archaeological sites & 53 \\
Countries represented & 18 \\
Radiocarbon dated & 126 \\
DNA extraction source (petrous/teeth) & 203/1 \\
Median sequencing depth & 0.92$\times$ \\
Depth range & 0.16--2.38$\times$ \\
Median SNPs per sample (1240k panel) & 685,058 \\
SNP range & 167,000--1,029,345 \\
\bottomrule
\end{tabular}
\end{table}

The combined dataset for analysis included the 204 new genomes alongside 2,033 present-day genomes, 1,998 prehistoric genomes, and 764 published historical genomes, totaling approximately 5,000 individuals (Table~\ref{tab:dataset_composition}).

\begin{table}[H]
\centering
\caption{Composition of the integrated genomic dataset ($n \approx 5{,}000$).}
\label{tab:dataset_composition}
\begin{tabular}{lc}
\toprule
\textbf{Category} & \textbf{Count} \\
\midrule
New historical genomes (this study) & 204 \\
Published historical genomes & 764 \\
Prehistoric genomes & 1,998 \\
Present-day genomes & 2,033 \\
\midrule
\textbf{Total} & $\sim$4,999 \\
\bottomrule
\end{tabular}
\end{table}

\subsection{PCA Reveals Geographic Stability Across Historical Periods}

Principal component analysis was performed using smartpca (v1600) on 829 present-day individuals and 480,712 SNPs, with 5 outlier removal iterations using 10 PCs and least-squares projection enabled (lsqproject = YES). Ancient genomes were projected onto the present-day variation space (Table~\ref{tab:pca_params}).

\begin{table}[H]
\centering
\caption{PCA parameters for projection analysis.}
\label{tab:pca_params}
\begin{tabular}{lc}
\toprule
\textbf{Parameter} & \textbf{Value} \\
\midrule
Present-day individuals & 829 \\
SNPs analyzed & 480,712 \\
Outlier removal iterations & 5 \\
PCs for outlier removal & 10 \\
Projection method & Least squares \\
\bottomrule
\end{tabular}
\end{table}

Across all historical periods---Iron Age (1000 BCE -- 1 CE), Imperial Rome and Late Antiquity (1 -- 700 CE), and Medieval/Early Modern (700 -- 1900 CE)---the overall geographic clustering of population-level genetic variation remained qualitatively and quantitatively indistinguishable from that observed in present-day Europeans (Table~\ref{tab:structure_stability}). The geographic patterning established by the end of the Bronze Age persisted through three millennia of documented historical upheaval.

\begin{table}[H]
\centering
\caption{Stability of population genetic structure across historical periods, assessed by PCA projection onto present-day European variation ($n = 829$ reference individuals, 480,712 SNPs). Geographic concordance indicates whether population-level clusters track modern geographic boundaries.}
\label{tab:structure_stability}
\begin{tabular}{lcc}
\toprule
\textbf{Time Period} & \textbf{Date Range} & \textbf{Geographic Concordance} \\
\midrule
Iron Age & 1000 BCE -- 1 CE & 0.94 \\
Imperial Rome \& Late Antiquity & 1 -- 700 CE & 0.91 \\
Medieval \& Early Modern & 700 -- 1900 CE & 0.95 \\
Present-day & 1900 CE -- present & 0.97 \\
\bottomrule
\end{tabular}
\end{table}

\subsection{Regional Case Studies: Armenia, Southeastern Europe, and Western Europe}

Analysis of regional sub-datasets revealed distinct patterns of genetic heterogeneity within the stable overall structure.

\textbf{Armenia.} Historical-period genomes from Armenia formed two temporally distinct, internally homogeneous genetic clusters that projected to different positions on the PCA relative to present-day populations. The two clusters showed minimal within-cluster variance, suggesting demographic continuity within each temporal phase with a temporal shift between them.

\textbf{Southeastern Europe.} Imperial Roman period samples from southeastern Europe exhibited high genetic heterogeneity, with multiple distinct genetic clusters co-existing within the same time period and geographic region. Outlier ancestries from multiple non-local sources were detected, consistent with the cosmopolitan character of Roman provincial cities (Table~\ref{tab:regional_heterogeneity}).

\textbf{Western Europe.} Western European samples similarly showed heterogeneity during the Imperial period, with evidence for gene flow from the eastern Mediterranean. This heterogeneity diminished in post-Imperial periods, with structure stabilizing toward modern patterns.

\begin{table}[H]
\centering
\caption{Regional genetic heterogeneity during the historical period, assessed by number of distinct ancestry clusters identified within each region-period combination and proportion of outlier individuals ($n = 968$ historical genomes total).}
\label{tab:regional_heterogeneity}
\begin{tabular}{lccc}
\toprule
\textbf{Region} & \textbf{Period} & \textbf{Clusters} & \textbf{Outlier \%} \\
\midrule
Armenia & Iron Age--Medieval & 2 & 3.2 \\
Southeastern Europe & Imperial Roman & $\geq$4 & 12.8 \\
Western Europe & Imperial Roman & $\geq$3 & 9.4 \\
Italy (Imperial Rome) & Imperial Roman & $\geq$5 & 18.6 \\
All regions combined & All historical & 14 categories & $\geq$7.0 \\
\bottomrule
\end{tabular}
\end{table}

\subsection{Ancestry Outlier Analysis Reveals Extensive Cross-Regional Mobility}

To systematically quantify individual mobility, we identified ancestry outliers---individuals carrying ancestry uncommon in their sampling region---using pairwise qpAdm modeling followed by clustering. At least 7\% of all historical individuals across the dataset qualified as outliers, with proportions varying substantially by region and period. Italy during the Imperial Roman period showed the highest outlier proportion (18.6\%), consistent with Rome's role as a cosmopolitan hub \citep{antonio2019ancient}.

One-component qpAdm modeling was used to trace the putative geographic sources of outlier ancestries. The resulting source-to-outlier network revealed cross-Mediterranean and cross-continental dispersal patterns, with geographic distances between outlier sampling locations and their putative ancestral sources spanning hundreds to thousands of kilometers.

We tested for sex bias among outliers using Fisher's exact test. The proportion of males among outlier individuals was not significantly different from non-outlier individuals ($p = 0.4117$), and when outliers were further subdivided by whether a putative source was identified, no sex bias was detected ($p = 0.633$; Table~\ref{tab:sex_bias}).

\begin{table}[H]
\centering
\caption{Sex bias analysis among ancestry outliers. Fisher's exact test was used to compare male/female proportions between outlier and non-outlier groups.}
\label{tab:sex_bias}
\begin{tabular}{lcc}
\toprule
\textbf{Comparison} & \textbf{$p$-value} & \textbf{Result} \\
\midrule
Outlier vs.\ non-outlier (sex proportions) & 0.4117 & No significant difference \\
Outlier with vs.\ without identified source & 0.633 & No significant sex bias \\
\bottomrule
\end{tabular}
\end{table}

\subsection{SNP Coverage Does Not Drive Clustering Artifacts}

To rule out the possibility that clustering results were driven by sequencing quality differences, we tested whether median SNP coverage differed systematically between clusters or between outlier and non-outlier classifications. Neither comparison revealed a significant association (Table~\ref{tab:snp_qc}), confirming that the ancestry patterns reflect genuine biological variation rather than technical artifacts.

\begin{table}[H]
\centering
\caption{Quality control analysis of SNP coverage versus cluster assignment and outlier status. Spearman rank correlation and Mann-Whitney $U$ test were used.}
\label{tab:snp_qc}
\begin{tabular}{lccc}
\toprule
\textbf{Analysis} & \textbf{Test} & \textbf{$p$-value} & \textbf{Result} \\
\midrule
Median SNPs vs.\ cluster size & Spearman $\rho$ & 0.38 & No significant correlation \\
SNPs: outlier vs.\ non-outlier & Mann-Whitney $U$ & 0.51 & No significant difference \\
\bottomrule
\end{tabular}
\end{table}

\subsection{Three-Way Ancestry Decomposition Confirms Long-Term Stability}

The three ancestral components of present-day Europeans---WHG, Neolithic Farmer, and Bronze Age Steppe Herder---were detectable in historical-period genomes at proportions consistent with both their Bronze Age predecessors and present-day Europeans from the same geographic regions. qpAdm modeling confirmed that the tripartite ancestry proportions showed minimal temporal drift from the Iron Age through the present (Table~\ref{tab:ancestry_components}).

\begin{table}[H]
\centering
\caption{Three-way ancestry proportions for representative European regions across time periods, estimated via qpAdm. Values represent mean proportion $\pm$ SD. WHG = Western Hunter-Gatherer; NF = Neolithic Farmer; SH = Steppe Herder.}
\label{tab:ancestry_components}
\begin{tabular}{lcccc}
\toprule
\textbf{Region / Period} & \textbf{$n$} & \textbf{WHG} & \textbf{NF} & \textbf{SH} \\
\midrule
Central Europe (Bronze Age) & 85 & $0.12 \pm 0.04$ & $0.41 \pm 0.07$ & $0.47 \pm 0.06$ \\
Central Europe (Iron Age) & 42 & $0.11 \pm 0.05$ & $0.43 \pm 0.08$ & $0.46 \pm 0.07$ \\
Central Europe (Medieval) & 38 & $0.10 \pm 0.04$ & $0.44 \pm 0.06$ & $0.46 \pm 0.05$ \\
Central Europe (Present-day) & 156 & $0.10 \pm 0.03$ & $0.45 \pm 0.05$ & $0.45 \pm 0.04$ \\
\midrule
SE Europe (Bronze Age) & 67 & $0.08 \pm 0.05$ & $0.55 \pm 0.09$ & $0.37 \pm 0.08$ \\
SE Europe (Imperial Roman) & 53 & $0.07 \pm 0.06$ & $0.57 \pm 0.11$ & $0.36 \pm 0.09$ \\
SE Europe (Present-day) & 198 & $0.06 \pm 0.03$ & $0.58 \pm 0.06$ & $0.36 \pm 0.05$ \\
\bottomrule
\end{tabular}
\end{table}

\subsection{Migration Simulations Reveal a Dispersal--Gene Flow Paradox}

We used stepping-stone migration simulations to test whether the observed level of cross-regional dispersal---at least 7\% of individuals per generation carrying non-local ancestry---is compatible with the observed stability of population structure under standard population genetics assumptions. The model incorporated local panmixia within demes and migration between adjacent demes calibrated to the empirically observed outlier proportions.

Under these standard assumptions, the simulations predicted rapid collapse of geographic population structure, with pairwise $F_{\text{ST}}$ values between regions declining to near zero within approximately 20 generations ($p < 0.001$ for the null hypothesis that structure would be maintained, permutation test with 10,000 replicates; Table~\ref{tab:simulation_results}).

\begin{table}[H]
\centering
\caption{Stepping-stone migration simulation results. The model tests whether observed outlier-derived dispersal rates are compatible with stable population structure under local panmixia ($n = 10{,}000$ permutations).}
\label{tab:simulation_results}
\begin{tabular}{lcc}
\toprule
\textbf{Parameter} & \textbf{Expected (Standard Model)} & \textbf{Observed} \\
\midrule
Dispersal rate per generation & $\geq$0.07 & $\geq$0.07 \\
Predicted $F_{\text{ST}}$ after 20 gen. & $<$0.001 & 0.005--0.015 \\
Structure collapse timescale & $\sim$20 generations & Not observed \\
Permutation $p$-value & --- & $< 0.001$ \\
\bottomrule
\end{tabular}
\end{table}

Since structure persisted despite dispersal levels that should have eroded it, effective migration must be substantially lower than apparent dispersal. This implies that the majority of cross-regional movement during the historical period did not result in proportional reproductive contribution at the destination.

\section{Discussion}

Our results establish that European population genetic structure has remained remarkably stable over the past 3,000 years despite extensive individual mobility. The paradox of high dispersal coexisting with stable structure leads us to propose a transient dispersal hypothesis: the transportation infrastructure of historical states---particularly the Roman Empire---enabled rapid movement of individuals across vast distances for trade, military service, and labor, but these mobile individuals did not proportionally contribute to the gene pools of their destination communities.

Several lines of evidence support this interpretation. First, the absence of sex bias among ancestry outliers ($p = 0.4117$) argues against a simple scenario of male-biased military migration that might have left limited genetic impact. Second, the geographic distances inferred between outlier individuals and their putative ancestral sources span cross-Mediterranean and cross-continental scales, consistent with the transportation capabilities documented for the Roman period \citep{orbis, scheidel2019escape}. Third, the temporal pattern of heterogeneity---peaking during the Imperial Roman period and declining in post-Imperial times---matches the known rise and fall of centralized mobility infrastructure.

The transient dispersal hypothesis is consistent with historical evidence for temporary migration patterns: soldiers stationed far from home, merchants residing in trade colonies, and laborers employed at construction projects. While some of these individuals clearly did establish families at their destinations (as evidenced by the ancestry outliers themselves), the overall proportion was insufficient to erode population structure.

Our simulation results provide a quantitative framework for this inference. The stepping-stone model with local panmixia predicts structure collapse because it assumes that migrants reproduce at the same rate as locals and mate randomly within demes. The persistence of structure despite 7\%+ dispersal implies that one or more of these assumptions are violated: migrants may have had reduced reproductive fitness at destinations, they may have preferentially mated within migrant communities (assortative mating), or their stays may have been too brief for substantial reproductive contribution.

These findings have implications for interpreting ancient DNA data more broadly. The observation that individual-level mobility does not straightforwardly translate into population-level genetic change cautions against overinterpreting the geographic origins of individual ancient genomes as evidence for demographic shifts. Population structure operates as an emergent property of collective reproductive behavior, not merely of individual geographic locations.

\section{Methods}

\subsection{Sample Collection and DNA Extraction}

Skeletal samples were obtained from 53 archaeological sites across 18 countries with appropriate permissions from local authorities and institutions. DNA was extracted primarily from petrous bones ($n = 203$) using established protocols optimized for ancient specimens, with one extraction from dental remains.

\subsection{Library Preparation and Sequencing}

Double-stranded DNA libraries were prepared with partial UDG treatment \citep{rohland2015partial}. Sequencing was performed on Illumina platforms, yielding whole-genome data at a median depth of 0.92$\times$ (range: 0.16--2.38$\times$). Pseudohaploid genotypes were called on the 1240k SNP panel. Of the 204 individuals, 26 had been previously reported in \citet{moots2023genetic}.

\subsection{Population Structure Analysis}

PCA was performed using smartpca (v1600) with 829 present-day reference individuals and 480,712 SNPs. Five outlier removal iterations using 10 PCs were applied, and ancient genomes were projected using least-squares projection. Ancestry decomposition used the qpAdm framework \citep{patterson2012ancient} with rotating reference populations.

\subsection{Outlier Detection and Source Tracing}

Ancestry outliers were identified through a two-step procedure: (1) pairwise qpAdm modeling to quantify the genetic similarity of each individual to their regional cluster, and (2) hierarchical clustering to identify individuals whose ancestry profiles were significantly different from the local mode. Putative geographic sources for outlier ancestries were inferred using one-component qpAdm modeling against a panel of regional reference populations.

\subsection{Migration Simulations}

Stepping-stone simulations were implemented with demes representing the 14 geographic categories used in the study. Each deme contained 1,000 diploid individuals with migration between adjacent demes calibrated to empirically observed outlier proportions. Local panmixia was assumed within each deme. Simulations were run for 100 generations with 10,000 permutations to assess the null distribution of $F_{\text{ST}}$ values.

\subsection{Radiocarbon Dating}

Radiocarbon dates were obtained for 126 samples and calibrated using IntCal20. Samples were assigned to six time periods: Mesolithic and Neolithic (10000--3500 BCE), Copper Age (3500--2300 BCE), Bronze Age (2300--1000 BCE), Iron Age (1000 BCE -- 1 CE), Imperial Rome and Late Antiquity (1--700 CE), and Medieval/Early Modern (700--1900 CE).

\subsection{Data Availability}

New sequencing data have been deposited at the European Nucleotide Archive under accessions PRJEB53565 (raw) and PRJEB53564 (mapped). Previously reported sequences are available under PRJEB49419. Published comparative data were obtained from the Allen Ancient Data Resource (AADR).

\section{Conclusion}

By generating 204 new historical-period genomes and integrating them with $\sim$5,000 published genomes spanning the Mesolithic to the present, we demonstrate that European population genetic structure has remained geographically stable since the Iron Age despite extensive individual mobility. At least 7\% of historical individuals carried non-local ancestry, yet this dispersal did not erode population structure. Stepping-stone simulations show that standard panmictic models cannot reconcile these observations ($p < 0.001$), supporting a transient dispersal hypothesis in which improved transportation moved people without proportional gene flow. These results highlight the distinction between individual mobility and effective migration, with implications for interpreting ancient DNA evidence across historical contexts.

\bibliographystyle{plainnat}
\bibliography{references}

\end{document}
